\documentclass[11pt,a4paper]{article}
\usepackage[margin=0.85in,top=1in,bottom=1in]{geometry}
\usepackage{amsmath,amssymb,amsfonts}
\usepackage{graphicx}
\usepackage{float}
\usepackage{enumitem}
\usepackage{microtype}
\usepackage{caption}
\usepackage[hidelinks]{hyperref}
\usepackage[most,breakable]{tcolorbox}
\tcbuselibrary{breakable,skins}
\usepackage{fontspec}
\setmainfont{DejaVu Serif}
\setsansfont{DejaVu Sans}
\setmonofont{DejaVu Sans Mono}
\captionsetup{font=small,labelfont=bf,skip=4pt}
\setlist{leftmargin=1.5em,topsep=2pt}

%% Breakable problem card — prevents content cutoff across pages
\newtcolorbox{solcard}[3]{
  breakable, enhanced,
  colback=white, colframe=gray!50, arc=2pt,
  title={\textbf{#1}\hfill\textsf{\small #3\,pts}},
  fonttitle=\normalsize\sffamily\bfseries,
  before upper={\small\textit{#2}\par\smallskip\hrule height 0.4pt\smallskip},
  top=4pt, bottom=6pt, left=7pt, right=7pt,
  parbox=false,
}

\title{\textbf{Oil production --- Solution}}
\author{\normalsize X24 (2024) — English translation}
\date{}
\begin{document}\maketitle\vspace{-0.5em}
\begin{solcard}{A1}{Let the oil reservoir be a section of ancient river sandstone deposits in the shape of a parallelepiped with a height of $h = 10~{\mathrm{m}}$, a width of $b = 100~{\mathrm{m}}$ and a length of $L = 2000~{\mathrm{m}}$. Rock porosity $\varphi = 0.1$. Estimate the oil reserves $m_{н}$ in this field. Express the answer in terms of $L$, $b$, $h$, $\rho$ and $\varphi$, and also give its numerical value in tons. Consider that the oil fluid completely fills the pore volume.}{0.30}

From the definition of porosity for the volume of oil we find: $$V_\text{н}=\varphi\cdot bhL{.} $$Since $m_\text{n}=\rho V_\text{n}$, we have:

\par\smallskip\noindent\textbf{Answer:} \textbackslash{}textbf{Answer:} $$m_\text{н}=\rho\varphi bhL=160\cdot 10^3~\text{тонн}{.} $$
\end{solcard}

\begin{solcard}{A2}{Let the reservoir oil pressure at the bottom of the deposits be $p_{пл}=250~{атм}$. Find at what maximum depth $H_{max}$ the deposit will be gushing, i.e. oil will flow to the surface under its own pressure. Express the answer in terms of $\rho$, $g$ and $p_{пл}$, and also give its numerical value. The compressibility of oil can be neglected.}{0.30}

The liquid will flow out under the influence of reservoir pressure if it exceeds the pressure of the hydrostatic column of oil: $$p_\text{пL}\geq \rho gh{.} $$Thus:

\par\smallskip\noindent\textbf{Answer:} \textbackslash{}textbf{Answer:} $$H_{max}=\cfrac{p_\text{пL}}{\rho g}\approx 3{.}125~\text{km} $$
\end{solcard}

\begin{solcard}{A3}{Estimate the maximum possible oil recovery factor $\alpha_{max}$ in the flowing mode at reservoir pressure $p_{пл}=250~{атм}$, if the oil compressibility is $\beta = 5\cdot 10^{-10}~{\mathrm{Pa}}$. Express the answer in terms of $\beta$ and $p_{пл}$, and also give its numerical value. Consider streambed sediments to be insulated by impermeable clays of low porosity. The depth of deposits $H$ can be chosen arbitrarily.}{0.60}

As already mentioned, the condition is that reservoir pressure is caused by fluid compression relative to normal conditions. If a fluid is under normal atmospheric pressure, its fixed amount of substance occupies the largest possible volume, which means that the largest proportion of displaced fluid is achieved if in the final state it is at the lowest possible pressure. This can be realized if the depth of deposits $H$ is directed to zero. We have: $$\alpha=\cfrac{dV}{V}\approx \beta\Delta{p}{.} $$Since the reservoir pressure is many times higher than atmospheric pressure, we find:

\par\smallskip\noindent\textbf{Answer:} \textbackslash{}textbf{Answer:} $$\alpha_{max}\approx \beta p_\text{пL}\approx 1{.}25\text{%}{.} $$
\end{solcard}

\begin{solcard}{A4}{Given the same data, estimate the maximum possible oil recovery factor $\alpha_{max}$ in the flowing mode, if there is water in the reservoir sediments below with a volume of $kV_0$ ($k = 9$) with initial oil reserves $V_0$. Consider the compressibility of water equal to the compressibility of oil. Express the answer in terms of $\beta$ and $p_{пл}$, and also give its numerical value. Consider that the liquid is taken from above, i.e. Only oil is taken. The depth of the deposits $H$ can be chosen arbitrarily.}{0.30}

Similar to the previous point, the depth of deposits $H$ should tend to zero. However, in this case, the oil recovery factor increases significantly, since not only the fluid, but also water is compressed, and since only oil is taken, its volume taken $dV$ consists of the change in the volumes of oil and water:: $$\alpha=\cfrac{dV}{V_0}=\cfrac{dV_\text{н}+dV_\text{в}}{V_0}=(1+9)\beta\Delta{p}{.} $$Thus:

\par\smallskip\noindent\textbf{Answer:} \textbackslash{}textbf{Answer:} $$\alpha_{max}=10\beta p_\text{пL}\approx 12{.}5\text{%}{.} $$
\end{solcard}

\begin{solcard}{B1}{Let us consider the horizontal flow of fluid along the $x$ axis between two parallel planes of height $h$. Distance between planes $w\ll{h}$. Determine the volumetric flow rate (hereinafter in all points of the problem - flow) of the liquid $Q$ through the cross section $wh$. Express your answer in terms of $\eta$, $w$, $h$ and the pressure gradient $dp(x)/dx$.}{1.00}

We obtain the distribution of fluid velocities along the coordinate $z$, directed perpendicular to the direction of flow from one plate to another. The beginning of the axis $z$ is located in the middle between the plates. From the condition of conservation of momentum we have: $$-hdzdp+\eta hdx\left(\cfrac{dv(z+dz)}{dz}-\cfrac{dv(z)}{dz}\right)=0{,} $$hence: $$\cfrac{d^2v}{dz^2}=\cfrac{1}{\eta}\cfrac{dp}{dx}{.} $$Integrating один times, we find: $$\cfrac{dv}{dz}=\cfrac{z}{\eta}\cfrac{dp}{dx}+A{.} $$Since velocity посередине между пластинами испытывает экстремум – $A=0$. Then повторно integrating, we find: $$v(z)=C+\cfrac{z^2}{2\eta}\cfrac{dp}{dx}{.} $$For/At $z=\pm w/2$ velocity должна обращаться в ноль, hence we find: $$v(z)=-\cfrac{1}{2\eta}\cfrac{dp}{dx}\left(\cfrac{w^2}{4}-z^2\right){.} $$For the flux $Q$ we have: $$Q=\int\limits_{-w/2}^{w/2}v(z)hdz=-\cfrac{h}{2\eta}\cfrac{dp}{dx}\int\limits_{-w/2}^{w/2}\left(\cfrac{w^2}{4}-z^2\right)dz{.} $$Integrating, we find:

\par\smallskip\noindent\textbf{Answer:} \textbackslash{}textbf{Answer:} $$Q=-\cfrac{w^3h}{12\eta}\cfrac{dp}{dx}{.} $$
\end{solcard}

\begin{solcard}{B2}{An excess pressure $\Delta p$ is created in the center of the gap. Find the dependence of the excess pressure $p'$ in the slot on the coordinate $x$. Express your answer using $\Delta{p}$, $Q$, $E$, $h$, $\eta$ and $x$.}{1.00}

Since the fluid flow spreads equally in two different directions, in each of them it is equal to $Q/2$. Let the $x$ axis originate at the center of the slit and be directed along one of the flows. Let's use the result of point $\mathrm{B1}$: $$\cfrac{Q}{2}=-\cfrac{w^3h}{12\eta}\cfrac{dp'}{dx}{.} $$Воспользовавшись эмпирическим соratioм for ширины щели, we get: $$\cfrac{Q}{2}=-\cfrac{h^4p'^3}{12\eta E^3}\cfrac{dp'}{dx}{.} $$Let us integrate поrayенное expression: $$\int\limits_{\Delta{p}}^{p'(x)}p'^3dp'=\cfrac{p'^4(x)-\Delta{p}^4}{4}=-\cfrac{6Q\eta E^3x}{h^4}{.} $$Thus:

\par\smallskip\noindent\textbf{Answer:} \textbackslash{}textbf{Answer:} $$p'(x)=\sqrt[4]{\Delta{p}^4-\cfrac{24Q\eta E^3x}{h^4}}{.} $$
\end{solcard}

\begin{solcard}{B3}{The crack ends at a position corresponding to zero overpressure. Determine the crack length $L$. Express your answer in terms of $\Delta{p}$, $E$, $h$, $\eta$ and $Q$.}{0.20}

At the maximum length of the crack, the pressure $p'$ becomes zero, whence the half-length of the crack $L/2$ is: $$\cfrac{L}{2}=\cfrac{\Delta{p}^4h^4}{24Q\eta E^3}{.} $$Thus:

\par\smallskip\noindent\textbf{Answer:} \textbackslash{}textbf{Answer:} $$L=\cfrac{\Delta{p}^4h^4}{12Q\eta E^3}{.} $$
\end{solcard}

\begin{solcard}{B4}{Determine the crack volume $V$. Express your answer in terms of $\Delta{p}$, $h$, $\eta$, $Q$ and $E$.}{0.70}

The volume of a crack is equal to twice the volume of its half: $$V=2\int\limits_{0}^{L}hw(x)dx=2\int\limits_{0}^{L}\cfrac{h^2p'}{E}dx{.} $$Let us use expressionм for $p'(x)$: $$V=\cfrac{2h^2\Delta{p}}{E}\int\limits_{0}^{L}\sqrt[4]{1-\cfrac{24Q\eta E^3x}{h^4\Delta{p}^4}}dx=\cfrac{h^6\Delta{p}^5}{12Q\eta E^4}\int\limits_{0}^{1}\sqrt[4]{1-z}dz{.} $$After elementary integration we find:

\par\smallskip\noindent\textbf{Answer:} \textbackslash{}textbf{Answer:} $$V=\cfrac{h^6\Delta{p}^5}{15Q\eta E^4}{.} $$
\end{solcard}

\begin{solcard}{B5}{Calculate the maximum possible values ​​of the crack length $L_{max}$ and its volume $V_{max}$.}{0.30}
\par\smallskip\noindent\textbf{Answer:} \textbackslash{}textbf{Answer:} $$L_{max}=\approx 83~\text{m}{.} $$
\end{solcard}

\begin{solcard}{С1}{Determine the speed $v$ of motion of the fluid boundary when the front moves by an amount $S$. Express your answer in terms of $p_1$, $p_2$, $L$, $\eta_1$, $\eta_2$, $k_1$ and $k_2$.}{1.00}

Darcy's law within the framework of the problem can be rewritten in the following form: $$v=-\cfrac{k}{\eta}\cfrac{\partial p}{\partial x}{.} $$Velocity движения $v$ одинакова for обеих жидкостей and постоянна по всей длине $L$. Then for градиентов давления в нефти and в воде taking into account equalsства их проницаемостей $k_1=k_2=k$ we get: $$\cfrac{\partial p_\text{н}}{\partial x}=-\cfrac{\eta_1v}{k_1}\qquad \cfrac{\partial p_\text{в}}{\partial x}=-\cfrac{\eta_2v}{k_2}{.} $$For полного изменения давления we find: $$p_2-p_1=\cfrac{\eta_2Sv}{k_2}+\cfrac{\eta_1(L-S)v}{k_1}{.} $$Thus:

\par\smallskip\noindent\textbf{Answer:} \textbackslash{}textbf{Answer:} $$v=\cfrac{p_2-p_1}{\cfrac{\eta_1L}{k_1}+\left(\cfrac{\eta_2}{k_2}-\cfrac{\eta_1}{k_1}\right)S}{.} $$
\end{solcard}

\begin{solcard}{C2}{Determine the dependence of the movement $S$ of the front on time $t$. Express your answer using $p_1$, $p_2$, $L$, $\eta_1$, $\eta_2$, $k$ and $t$}{0.90}

Since $v=dx/dt$ and $k_1=k_2=k$, we have: $$dt=\cfrac{(\eta_1L+(\eta_2-\eta_1)S)dS}{k(p_2-p_1)}{.} $$Integrating, let us find: $$t=\cfrac{1}{k(p_2-p_1)}\int\limits_{0}^S(\eta_1L+(\eta_2-\eta_1)x)dx=\cfrac{1}{k(p_2-p_1)}\left(\eta_1LS+\cfrac{(\eta_2-\eta_1)S^2}{2}\right){.} $$For/Atведём квадратное equation к классическому виду: $$S^2-\cfrac{2\eta_1LS}{\eta_1-\eta_2}+\cfrac{2k(p_2-p_1)t}{\eta_1-\eta_2}=0{.} $$Решая, we get: $$S(t)=\cfrac{\eta_1L}{\eta_1-\eta_2}\pm\sqrt{\left(\cfrac{\eta_1L}{\eta_1-\eta_2}\right)^2-\cfrac{2k(p_2-p_1)t}{\eta_1-\eta_2}}{.} $$Since $x(0)=0$ – let’s choose the root with a minus sign. Finally:

\par\smallskip\noindent\textbf{Answer:} \textbackslash{}textbf{Answer:} $$S(t)=\cfrac{\eta_1L}{\eta_1-\eta_2}-\sqrt{\left(\cfrac{\eta_1L}{\eta_1-\eta_2}\right)^2-\cfrac{2k(p_2-p_1)t}{\eta_1-\eta_2}}{.} $$
\end{solcard}

\begin{solcard}{C3}{Determine the total time $\tau$ for oil displacement from the field. Express the answer in terms of $p_1$, $p_2$, $L$, $\eta_1$, $\eta_2$ and $k$ and calculate it.}{0.50}

The total time $\tau$ of oil displacement is $t(L)$. Thus:

\par\smallskip\noindent\textbf{Answer:} \textbackslash{}textbf{Answer:} $$\tau=\cfrac{(\eta_1+\eta_2)L^2}{2k(p_2-p_1)}\approx 26~\text{Lет}{.} $$
\end{solcard}

\begin{solcard}{С4}{Under what condition on the parameters of the system will the motion of the boundary be stable, that is, if the shape of the boundary deviates small from flat, this deviation will not increase? Write the stability condition in terms of $\eta_1$, $\eta_2$, $k_1$ and $k_2$. Is the fluid flow considered in paragraphs $\mathrm{C2}$ and $\mathrm{C3}$ stable?}{0.80}

The deviation will not increase if $v(x+dx){<}v(x)$. Note that the numerator of the expression for speed is constant, and the denominator is a linear function of $x$ with slope $a=\eta_2/k_2-\eta_1/k_1$. The magnitude of the speed decreases with increasing $x$, if the value $a$ is positive, i.e., provided:

\par\smallskip\noindent\textbf{Answer:} \textbackslash{}textbf{Answer:} $$\cfrac{\eta_2}{k_2}{>}\cfrac{\eta_1}{k_1}{.} $$
\end{solcard}

\begin{solcard}{D1}{Find the dependence of the fluid flow velocity in such a pipe on the distance to the pipe axis $v(r)$, the maximum value of the velocity $v_{max}$ and the total fluid flow $Q$ through the cross section of the cylinder. Express your answers using $\Delta{p}$, $\eta$, $L$, $R$ and $r$.}{0.80}

Let us select a cylinder with radius $r$ and length $L$. Let us write down the condition for the constancy of its momentum: $$\pi r^2\Delta{p}+2\pi rL\eta\cfrac{\partial v}{\partial r}=0\Rightarrow \cfrac{\partial v}{\partial r}=-\cfrac{r}{2}\cfrac{\Delta{p}}{\eta L}{.} $$Integrating, we get: $$v=C-\cfrac{\Delta{p}r^2}{4\eta L}{.} $$Since $v(R)=0$, we obtain:


For flow $Q$ we have: $$Q=\int\limits_0^Rv(r)\cdot 2\pi rdr=\cfrac{\pi\Delta{p}}{2\eta L}\int\limits_{0}^R (R^2-r^2)rdr{.} $$Integrating, we find:

\par\smallskip\noindent\textbf{Answer:} \textbackslash{}textbf{Answer:} $$v(r)=\cfrac{\Delta{p}(R^2-r^2)}{4\eta L}{.} $$
\end{solcard}

\begin{solcard}{D2}{Express the distribution of fluid velocity $v(r)$ in terms of the total flow $Q$, $R$ and $r$.}{0.20}

We express the combination $\Delta{p}/(4\eta L)$ in terms of the full flow $Q$: $$\cfrac{\Delta{p}}{4\eta L}=\cfrac{2Q}{\pi R^4}{.} $$Substituting в expression for $v(r)$, we find:

\par\smallskip\noindent\textbf{Answer:} \textbackslash{}textbf{Answer:} $$v(r)=\cfrac{2Q}{\pi R^2}\left(1-\cfrac{r^2}{R^2}\right){.} $$
\end{solcard}

\begin{solcard}{D3}{Find the flow $Q$ in the face section at a distance $h$ from its lower edge and the corresponding expression for the vertical velocity $v(r,h)$ depending on the distance to the axis $r$ and the height $h$. Express your answers using $Q_0$, $H$, $R$, $r$ and $h$.}{0.20}

Since the liquid enters the cylinder uniformly along its side surface, we have: $$\cfrac{dQ}{dh}=\cfrac{Q_0}{H}{.} $$Since $Q(0)=0$, we find:


Let's use the expression obtained in step $\mathrm{D2}$: $$v(r{,}h)=\cfrac{2Q(h)}{\pi R^2}\left(1-\cfrac{r^2}{R^2}\right){.} $$Substituting $Q(h)$, we get:

\par\smallskip\noindent\textbf{Answer:} \textbackslash{}textbf{Answer:} $$Q(h)=\cfrac{Q_0h}{H}{.} $$
\end{solcard}

\begin{solcard}{D4}{Consider a ring of height $dh$ with inner and outer radii $r$ and $r+dr$ respectively. Using the fact that the fluid is incompressible, show that from the condition that the volume of the fluid inside the selected ring is constant, the following relation follows: $$\cfrac{\partial v}{\partial h}=-\cfrac{1}{r}\cfrac{\partial (u_rr)}{\partial r}{.}$$You can use this relation even if you could not prove it.}{0.30}

Let us write the expression for the total fluid flow $q$ into the ring: $$q=2\pi rdrv(r{,}h+dh)-2\pi rdrv(r{,}h)+2\pi dh(r+dr)u_r(r+dr{,}h)-2\pi dh ru_r(r{,}h)=0{.} $$Thus: $$2\pi rdrdv+2\pi drd(ru_r)=0{,} $$or: $$\cfrac{\partial v}{\partial h}+\cfrac{1}{r}\cfrac{\partial (ru_r)}{\partial r}=0{.} $$Which is what we needed to prove.

\end{solcard}

\begin{solcard}{D5}{Find the radial velocity of the fluid flow $u_r(r,h)$ depending on the distance to the axis $r$ and the height $h$, as well as the maximum value of its modulus $u_{r(max)}$. Express your answers using $Q_0$, $R$, $H$, $h$ and $r$.}{0.50}

Consistently use the results of points $\mathrm{D4}$ and $\mathrm{D3}$: $$\cfrac{\partial (ru_r)}{\partial r}=-r\cfrac{\partial v}{\partial h}=-\cfrac{2Q_0r}{\pi R^2H}\left(1-\cfrac{r^2}{R^2}\right){.} $$Note that for/at $r=0$ quantity $ru_r$ also equals нулю. Временно переdenote $r$ на $z$ в подынтегральной функции and we get: $$ru_r(r{,}h)=-\cfrac{2Q_0}{\pi R^2H}\int\limits_{0}^r\left(1-\cfrac{z^2}{R^2}\right)zdz=-\cfrac{Q_0}{\pi R^2H}\left(r^2-\cfrac{r^4}{2R^2}\right){.} $$Thus:


The maximum value of $u_r$ is achieved when its derivative with respect to $r$ is equal to zero: $$\cfrac{du_r}{dr}=0\Rightarrow 1-\cfrac{3r^2}{2R^2}=0\Rightarrow r_{max}=\sqrt{\cfrac{2}{3}}R{.} $$Then for максимального значения радиальной компоненты скорости $u_{max}$ we find:

\par\smallskip\noindent\textbf{Answer:} \textbackslash{}textbf{Answer:} $$u_r(r{,}h)=-\cfrac{Q_0}{\pi R^2H}\left(r-\cfrac{r^3}{2R^2}\right){.} $$
\end{solcard}

\begin{solcard}{D6}{What is the ratio $u_{r(max)}/v_{max}$? Express your answer in terms of $R$ and $H$.}{0.10}
\par\smallskip\noindent\textbf{Answer:} \textbackslash{}textbf{Answer:} $$\cfrac{u_{r(max)}}{v_{max}}=\cfrac{\sqrt{2}R}{3\sqrt{3}H} $$
\end{solcard}

\end{document}